% NOETHER PAPER 4 — KOREAN TRANSLATION-PRODUCER DRAFT — UNCHECKED
% Unit: T03-U10; current-authority whole lines 3780--3787
% Source slice: 2,560 LF UTF-8 bytes; SHA-256 39F6BBAC7FA2466F8D86F5EBDE759CEBA517AA29F4D26F88386293625AFBC6CA
% Pointer: NOETH-DE-AUTH-v004-20260804; SHA-256 A1C62FDACAA34DFC1B806DC18258F2E732539F7AE9D85AA4BD9E1067B8749D9F
% This is producer translation only: no source/Korean/formula review, compilation, or rendering.

\section*{\S\ 3. 기호 항등식\footnote{문제의 정식화와 문헌에 대해서는 서론을 참조하라.}.}
이제 기호적 방법의 둘째 기본정리를 옮겨야 한다. 즉, 모든 열 $S^\rho,p_\rho$를 분해 불가능한 것으로 보았을 때의 기약인 분해 불가능 항등식 전체를 세워야 한다. 이를 위하여 다음을 추가로 정한다. 기호열과 변수열의 의미에 맞추어, 변수열 $p_\rho$를 $p_{\rho-1}y$로 특수화한 뒤 생기는 식에 이 계의 항등식 가운데 하나를 적용하여 얻는 항등식은 합성된 것으로 본다. 반면 $k$개의 기호열을 포함하는 항등식은, 위 과정으로 더 적은 기호열을 가진 항등식에서 생기더라도 분해 불가능한 것으로 본다. 열 $S^\rho,p_\rho$가 항등식에 반드시 명시적으로 들어갈 필요는 없다. 나타나는 행렬곱들이 이 열들에만 의존한다는 것을 보이면, 이들을 분해 불가능한 양으로 이해하기에 충분하다. 그러나 우리는 이 경우에도, 일부에는 대입 (11.)을 적용하여, 언제나 명시적 표현을 얻을 수 있음을 보일 것이다. 일반 항등식은 Pascal이 제시한 특수 항등식들로부터 행렬 표현의 원리를 써서 얻는다. 이 특수 항등식들은 열 $x$와 $u$만을 포함하며, Study가 삼항 영역에서 제시한 항등식을 직접 옮긴 것이다.
\begin{align}
\sum_{i=1}^n(-1)^{i+1}\pair{a^{(i)}}{x}(a^{(1)}\cdots a^{(i-1)}a^{(i+1)}\cdots a^{(n)}u)&=(-1)^{n+1}\pair{u}{x}(a^{(1)}\cdots a^{(n)}),\tag{13}\\
\sum\pm\pair{a^{(1)}}{x^{(1)}}\pair{a^{(2)}}{x^{(2)}}\cdots\pair{a^{(n)}}{x^{(n)}}&=(a^{(1)}\cdots a^{(n)})(x^{(1)}\cdots x^{(n)}),\tag{14}\\
\sum_{i=1}^n(-1)^{i+1}\pair{u}{a^{(i)}}(a^{(1)}\cdots a^{(i-1)}a^{(i+1)}\cdots a^{(n)}x)&=(-1)^{n+1}\pair{u}{x}(a^{(1)}\cdots a^{(n)}).\tag{15}
\end{align}
(13.)과 (15.)에 각각 $(a^{(i)}\mid x)=(a^{(i)}q_{n-1})$와 $(u\mid a^{(i)})=(p_{n-1}a^{(i)})$를 대입하면 두 항등식을 더 얻지만, 우리의 관점에서는 본질적으로 다른 것으로 보지 않는다. (14.)는 행렬식의 곱정리를 나타낸다. 이제 (13.)과 (14.)를(쌍대적으로는 (15.)와 (14.)를) 적절히 만든 행렬의 곱정리 계로 어떻게 바꿀 수 있는지 보이겠다. 곱정리를 한 번은 $\rho$열 행렬에, 한 번은 $(\rho-1)$열 행렬에 적용하면 다음을 얻는다.
